% ===================== INTRODUCCION =====================
\section{Introduccion}
\label{sec:intro}

[La introduccion debe situar al lector en el contexto tecnico e industrial del proyecto.
Incluir: (1) descripcion del sistema o proceso objeto de estudio, (2) antecedentes
relevantes y referencias a estudios previos, (3) motivacion del estudio actual
(problema a resolver, cambio tecnologico, requerimiento normativo), y (4) estructura
del documento indicando que contiene cada capitulo.]

% ─────────────────────────────────────────────────────────────────────────────
\subsection{Antecedentes}

[Describir el historial del sistema, diagnosticos previos, modificaciones recientes
o cualquier informacion de referencia que contextualice el presente estudio.
Citar referencias bibliograficas con cite{clave-de-referencia}.]

% ─────────────────────────────────────────────────────────────────────────────
\subsection{Descripcion del Sistema}

[Describir el sistema de proceso, sus componentes principales, modo de operacion
y variables criticas. Incluir diagrama de flujo o P\&ID si aplica.]

% ─────────────────────────────────────────────────────────────────────────────
\subsection{Motivacion del Estudio}

[Explicar el evento, requerimiento o decision que origino la necesidad del presente
estudio. Identificar claramente la pregunta tecnica que el informe responde.]

% ─────────────────────────────────────────────────────────────────────────────
\subsection{Estructura del Documento}

[Describir brevemente el contenido de cada capitulo del informe para guiar al lector.]
